% 論文節の日本語訳（読み物 PDF）用 preamble
% lualatex + jlreq
\documentclass[12pt,paper=a4]{jlreq}
\usepackage{luatexja-fontspec}
\usepackage{amsmath,amssymb,bm}
\usepackage{graphicx}
\usepackage{hyperref}
\usepackage{enumitem}
\usepackage{booktabs}
\usepackage{microtype}
\usepackage{xcolor}
\usepackage{caption}
\usepackage{tcolorbox}
\tcbuselibrary{breakable,skins}

\hypersetup{
  colorlinks=true,
  linkcolor=blue!60!black,
  urlcolor=blue!60!black,
  citecolor=blue!60!black,
}

\graphicspath{{images/}{./}}

\setlength{\parindent}{1em}
\setlength{\parskip}{0.55em}
\setlist[itemize]{leftmargin=1.6em,itemsep=0.25em,topsep=0.35em}
\setlist[enumerate]{leftmargin=1.8em,itemsep=0.25em,topsep=0.35em}

\captionsetup{
  font=small,
  labelfont=bf,
  skip=0.4em,
  justification=raggedright,
  singlelinecheck=false,
}

\newcommand{\vx}{\bm{x}}
\newcommand{\vu}{\bm{u}}
\newcommand{\veps}{\bm{\epsilon}}
\newcommand{\vw}{\bm{w}}

\DeclareMathOperator*{\argmax}{arg\,max}
\DeclareMathOperator*{\argmin}{arg\,min}
\DeclareMathOperator{\Tr}{Tr}
\DeclareMathOperator{\diag}{diag}

% 命題・定理ボックス（ページ跨ぎ時も切れ目を抑える）
\newtcolorbox{propositionbox}[1][]{%
  enhanced,
  breakable,
  colback=blue!4!white,
  colframe=blue!45!black,
  coltitle=white,
  fonttitle=\bfseries\small,
  title={#1},
  attach boxed title to top left={yshift=-2mm, xshift=3mm},
  boxed title style={
    colback=blue!45!black,
    colframe=blue!45!black,
    boxrule=0pt,
    arc=1pt,
  },
  boxrule=0.5pt,
  arc=1.5pt,
  left=5pt, right=5pt, top=8pt, bottom=5pt,
  before skip=0.9em,
  after skip=0.9em,
  pad at break=3mm,
}

% 図はフロートせず本文位置に置く（分断を防ぐ）
\newcommand{\translationfigure}[2]{%
  \begin{center}
    \includegraphics[width=0.88\linewidth]{#1}\\[0.35em]
    {\small #2\par}
  \end{center}
  \vspace{0.4em}%
}

\begin{document}
{\color{gray}\footnotesize 出典: §4.5\par}
\vspace{0.2em}

\section*{4.5 蓄積摂動核を通じた直接的見方}


DDPM と Score SDE は、時間局所的にノイズがどう加えられるかを指定することで導入されるのが一般的である：DDPM では離散前向き遷移 $p(x_t\mid x_{t-\Delta t})$ を通じて、Score SDE では連続時間 SDE を通じて。しかし学習においては、とくに有用な対象は蓄積摂動核

\[
p_t(x_t\mid x_0)
\]


であり、これは元データサンプルで直接条件付けられた時刻 $t$ における分布を記述する。これはまさに、DDPM および連続時間デノイジング目的の式 (2.2.8) と式 (4.3.1) に自然に現れる形である。

これは補完的な視点を示唆する：局所ノイズ注入から始めてから $p_t(x_t\mid x_0)$ を導出するのではなく、蓄積摂動核を直接指定できる。以下で考える線形–ガウス過程については、許容可能な蓄積核が対応する線形 SDE を定め、その逆も成り立つ。この視点は各時刻における信号とノイズの量を明示にし、実践的学習目的に直接つながり、また第5章の Flow Matching の確率経路視点への有用な橋渡しも提供する。

\translationfigure{fig-4-8.jpg}{図4.8: 同じアフィン前向き過程の二つの見方。許容可能なアフィン拡散スケジュールについて、連続時間線形 SDE を通じた増分ノイズ注入と、式 (4.5.1) における蓄積摂動核の直接指定は、同じ条件付き分布 $p_t(x_t\mid x_0)$ を記述する。}


**一般アフィン前向き過程 $p_t(x_t\mid x_0)$。** まず蓄積前向き摂動核を定義することから始める：

\[
p_t(x_t\mid x_0):=\mathcal{N}\bigl(x_t;\, \alpha_t x_0,\, \sigma_t^2 I\bigr), \quad (4.5.1)
\]


ここで $x_0\sim p_{\mathrm{data}}$ であり、$\alpha_t$ と $\sigma_t$ は $t\in[0,T]$ の非負スカラー関数である。標準的な前向き拡散設定では、

\[
\alpha_0=1, \qquad \sigma_0=0
\]


をとり、過程が正確にデータから始まるようにし、対応する比が用いられるときは $t\in(0,T]$ で $\alpha_t>0$ および $\sigma_t>0$ と仮定する。

等価に、式 (4.5.1) からのサンプルは次のように書ける：

\[
x_t=\alpha_t x_0 + \sigma_t\epsilon, \qquad \epsilon\sim\mathcal{N}(0,I).
\]


この単純な形は、いくつかのなじみ深い確率経路を含む：

\begin{itemize}
\item \textbf{VE（NCSN）}: $\alpha_t\equiv 1$、ノイズスケール $\sigma_t$ は時間とともに増加；
\item \textbf{VP（DDPM）}: $\alpha_t^2+\sigma_t^2=1$；
\item \textbf{線形ガウス補間}: $t\in[0,1]$ で $\alpha_t=1-t$、$\sigma_t=t$。データからノイズへの時間慣例で書かれたとき、Flow Matching および Rectified Flow で用いられるなじみ深い線形経路に対応する。
\end{itemize}


\subsection*{4.5.1 Score SDE への接続}


Score SDE については、ドリフトと拡散項から始めてモーメントの ODE を解く（第4.4.3節を参照）よりも、まず蓄積摂動核 $p_t(x_t\mid x_0)$ を指定し、次いで対応する SDE 係数を回復する方がしばしば直感的である。

許容可能なアフィンスケジュールについて、対応する前向き SDE は式 (4.4.2) における $x$ について線形な形をとる：

\[
dx(t)=\underbrace{f(t)x(t)}_{f(x(t),t)}\, dt + g(t)\, dw(t),
\]


ここで $f,g:[0,T]\to\mathbb{R}$ は時間の実数値関数である。

係数 $f(t)$ と $g(t)$ は $\alpha_t$ と $\sigma_t$ から解析的に回復できる。重要なことに、$\alpha_t$ と $\sigma_t$ のすべての選択が有効な拡散を定義するわけではない：含意される拡散分散 $g^2(t)$ は非負のままでなければならない。等価に、この線形–ガウス設定では、信号対雑音比 $\alpha_t^2/\sigma_t^2$ は時間とともに増加できない。次の補題がこの関係を精密にする。

\begin{propositionbox}[補題 4.5.1: 前向き摂動核 $\Leftrightarrow$ 線形 SDE]
$\alpha_0=1$、$\sigma_0=0$、および $\alpha_t>0$、$\sigma_t>0$ が $t\in(0,T]$ で微分可能であると仮定する。$\lambda_t:=\log\dfrac{\alpha_t}{\sigma_t}$ と定義する。

前向き摂動核

\[
x_t=\alpha_t x_0+\sigma_t\epsilon, \qquad \epsilon\sim\mathcal{N}(0,I)
\]

が与えられたとき、同じ条件付き核 $p_t(x_t\mid x_0)=\mathcal{N}(x_t;\,\alpha_t x_0,\,\sigma_t^2 I)$ をもつ対応する線形 SDE

\[
dx(t)=f(t)x(t)\, dt + g(t)\, dw(t)
\]

が存在する必要十分条件は、係数

\[
\begin{aligned}
f(t)&=\frac{d}{dt}\log\alpha_t, \\
g^2(t)&=\frac{d\sigma_t^2}{dt}-2\frac{d}{dt}\log\alpha_t\,\sigma_t^2
=-2\sigma_t^2\frac{d}{dt}\lambda_t \quad (4.5.2)
\end{aligned}
\]

が $t\in(0,T]$ で $g^2(t)\ge 0$ を満たすことである。等価に、$\lambda_t$ は非増加でなければならない。

逆に、線形 SDE が $\alpha_t>0$、$\sigma_t>0$ をもつ条件付き核 $\mathcal{N}(x_t;\,\alpha_t x_0,\,\sigma_t^2 I)$ をもつなら、その係数は必然的に式 (4.5.2) を満たす。
\end{propositionbox}


\textbf{補題の証明。}
第4.4.3節より、結果は平均および共分散 ODE $m'(t)=f(t)m(t)$ および $\mathbf{P}'(t)=2f(t)\mathbf{P}(t)+g^2(t)I$ を、$m(t)=\alpha_t x_0$ および $\mathbf{P}(t)=\sigma_t^2 I$ と整合させることで従う。■

\textbf{注。}
終端条件付き分布を $x_0$ に独立にし、規定されたガウス事前に等しくするには、$\alpha_T=0$ および $\sigma_T^2$ を事前分散に等しくする必要がある。
\[
\alpha_t=\exp\Bigl(\int_0^t f(u)\, du\Bigr)
\]

であるため、有限終端時刻で $\alpha_T=0$ を達成するには
\[
\int_0^T f(u)\, du=-\infty
\]

が必要である。したがって、ドリフトが誘導する収縮は終端時刻近傍で特異にならなければならず；とくに $f$ は $[0,T]$ 上で有界であり続けられない。

拡散係数は規定分散を通じて
\[
g^2(t)=(\sigma_t^2)'-2\frac{\alpha_t'}{\alpha_t}\sigma_t^2
\]

と関係する。

$\sigma_t^2$ が有限の非ゼロ終端分散に近づく通常の正則スケジュールでは、$(\sigma_t^2)'$ は有界のままであり、$\alpha_t'/\alpha_t\to-\infty$ であるため、これもまた

\[
g^2(t)\to\infty \qquad (t\to T)
\]


を与える。

より一般に、$g(t)$ の精密な振る舞いは選ばれたスケジュールに依存する。

$f$ と $g$ が有限区間 $[0,T]$ 上で有界のままなら、$\alpha_T>0$ であり、有限 $T$ で厳密な忘却は起きえない。実際には、したがって $x_0$ への残余依存が無視でき $p_T\approx p_{\mathrm{prior}}$ となる終端時刻を選ぶか、$\alpha_t\to 0$ が漸近的に起きる無限時間定式化を考える。

上の補題より、係数 $f(t)$ と $g(t)$ をもつ線形 SDE を通じて増分ノイズ注入を指定することは、パラメータ $\alpha_t$ と $\sigma_t$ で摂動核を定義することと数学的に等価である。拡散モデル文献では、これら二つの視点は互換的に用いられる。したがって次を結論する：

\begin{propositionbox}[観察 4.5.1:]
有効なアフィン拡散スケジュールについて、二つの記述は等価である：$\alpha_t$ と $\sigma_t$ を通じて蓄積摂動核 $p_t(x_t\mid x_0)$ を指定しても、$f(t)$ と $g(t)$ を通じて対応する線形 SDE を指定してもよい。
\end{propositionbox}


\subsection*{4.5.2 変分ベース拡散モデルへの接続}


ベイズの法則を通じて導かれた、DDPM からの核心的恒等式を再訪する：

\[
p(x_{t-\Delta t}\mid x_t,x)
= p(x_t\mid x_{t-\Delta t})\cdot\frac{p_{t-\Delta t}(x_{t-\Delta t}\mid x)}{p_t(x_t\mid x)}, \quad (4.5.3)
\]


任意の $x$（通常 $x\sim p_{\mathrm{data}}$）について。この逆条件付き $p(x_{t-\Delta t}\mid x_t,x)$ はモデリングの中心であり、扱いやすい学習ターゲットと効率的サンプリングの双方を可能にする。

DDPM は典型的に増分核 $p(x_t\mid x_{t-\Delta t})$ をまず定義するが、蓄積遷移 $p_t(x_t\mid x_0)$ はしばしば、とくに事前と損失設計について、より解釈可能で実践的な定式化を提供する。

\textbf{遷移核の導出。} 今これを連続時間設定へ拡張する。$0\le t<s\le T$ を二つの（連続）時間点とする。摂動核 $p_t(x_t\mid x_0)$ が与えられると、中間として前向き核 $p(x_s\mid x_t)$ を用い、式 (4.5.3)\footnote{この恒等式は、$s$ を一般のより早い時刻として扱うことで、自然に連続時間へ拡張する。} を適用することで任意の $x$ について逆条件付き $p(x_t\mid x_s,x)$ を計算できる。次の補題がこの導出を要約し、$\alpha_t^2+\sigma_t^2=1$ を仮定することなく補題 2.2.2 を拡張する。

\begin{propositionbox}[補題 4.5.2: 逆条件付き遷移核]
$0\le t<s\le T$ とする。逆条件付き遷移核は：

\[
p(x_t\mid x_s,x)=\mathcal{N}\bigl(x_t;\, \mu(x_s,x;s,t),\, \sigma^2(s,t)I\bigr),
\]

ここで

\[
\mu(x_s,x;s,t)
:=\frac{\alpha_{s|t}\sigma_t^2}{\sigma_s^2}\, x_s
+ \frac{\alpha_t\sigma_{s|t}^2}{\sigma_s^2}\, x, \qquad
\sigma^2(s,t):=\sigma_{s|t}^2\frac{\sigma_t^2}{\sigma_s^2}. \quad (4.5.4)
\]

ここで $\alpha_{s|t}$ と $\sigma_{s|t}$ は次で定義される：

\[
\alpha_{s|t}:=\frac{\alpha_s}{\alpha_t}, \qquad
\sigma_{s|t}^2:=\sigma_s^2-\alpha_{s|t}^2\sigma_t^2.
\]
\end{propositionbox}


\textbf{補題の証明。}
まず前向き遷移核を計算する：
\[
p(x_s\mid x_t)=\mathcal{N}\bigl(x_s;\, \alpha_{s|t}x_t,\, \sigma_{s|t}^2 I\bigr). \quad (4.5.5)
\]

逆核は次いでベイズの法則から従い、関与するすべての分布がガウスであるため、結果は直接計算で導出できる。さらなる詳細は (Kingma et al., 2021) の付録 A を参照。■

$p(x_{t+\Delta t}\mid x_t)$ と $p_t(x_t\mid x_0)$ は理論的に等価であるが、$p_t(x_t\mid x_0)$ がしばしばより中心的役割を果たす。式 (4.5.5) のステップごとの遷移は、主に閉形式の逆核を得るために機能する。近年の研究（Kingma et al., 2021）はしたがって、その明確さと解釈可能性のため $p_t(x_t\mid x_0)$ を直接指定することを好む。

\textbf{逆過程のモデリング、学習、およびサンプリング。} 学習目的（式 (2.2.13) の ELBO）および第2.2節で導入されたモデリング枠組みは、我々の一般化設定の下でも適用可能である。明確さのため、Kingma et al. (2021) に従い、$x_\phi(x_s,s)$ で表される $x$-予測定式化を採用する。しかし、関係（式 (2.2.12) と同様）

\[
x_s=\alpha_s x_\phi(x_s,s)+\sigma_s\epsilon_\phi(x_s,s), \qquad\text{任意の与えられた }x_s\sim q_s
\]


により、$\epsilon_\phi(x_s,s)$ で表される等価な $\epsilon$-予測視点もまた有効である。

**モデリングと拡散損失 $\mathcal{L}_{\mathrm{diffusion}}$。** DDPM と同様、式 (4.5.4) の条件付き分布 $p(x_t\mid x_s,x)$ は、クリーン信号 $x$ を学習可能な予測器 $x_\phi(x_s,s)$ で置き換えることを動機づけ、次の形のパラメータ化された逆モデルを与える：

\[
p_\phi(x_t\mid x_s):=\mathcal{N}\bigl(x_t;\, \mu_\phi(x_s,s,t),\, \sigma^2(s,t)I\bigr), \quad (4.5.6)
\]


平均は次のようにパラメータ化される：
\[
\mu_\phi(x_s,s,t)
=\frac{\alpha_{s|t}\sigma_t^2}{\sigma_s^2}\, x_s
+ \frac{\alpha_t\sigma_{s|t}^2}{\sigma_s^2}\, x_\phi(x_s,s).
\]


式 (4.5.1) の前向き核が与えられると、$\mathcal{L}_{\mathrm{diffusion}}(x;\phi)$ における KL ダイバージェンスは重み付き回帰損失に帰着する：

\begin{align*}
\begin{aligned}
D_{\mathrm{KL}}\bigl(p(x_t\mid x_s,x_0)\Vert p_\phi(x_t\mid x_s)\bigr)
&= \frac{1}{2\sigma^2(s,t)}\bigl\|\mu(x_s,x_0;s,t)-\mu_\phi(x_s,s,t)\bigr\|_2^2 \\
&= \tfrac12\bigl(\mathrm{SNR}(t)-\mathrm{SNR}(s)\bigr)\bigl\|x_0-x_\phi(x_s,s)\bigr\|_2^2, \quad (4.5.7)
\end{aligned}
\end{align*}


ここで $x_s=\alpha_s x_0+\sigma_s\epsilon$、$x_0\sim p_{\mathrm{data}}$、$\epsilon\sim\mathcal{N}(0,I)$、および $\mathrm{SNR}(s):=\alpha_s^2/\sigma_s^2$ は時刻 $s$ における信号対雑音比を表す。

\textbf{注。}
(Kingma et al., 2021) では、著者らは $t\to s$ における式 (4.5.7) の連続時間極限を研究し、次を得る：
\[
\mathcal{L}_{\mathrm{VDM}}^{\infty}(x_0)
= -\tfrac12\,\mathbb{E}_{s,\epsilon\sim\mathcal{N}(0,I)}\mathrm{SNR}'(s)\,\|x_0-x_\phi(x_s,s)\|_2^2.
\]

この設定はまた学習可能なノイズスケジュールを導入し、連続データを超えて一般化するが、そのような拡張は現在の議論の範囲外である。

\textbf{サンプリング。} サンプリングは式 (4.5.6) のパラメータ化された核を用い、DDPM と同様に進む：

\[
x_t
= \underbrace{
\frac{\alpha_{s|t}\sigma_t^2}{\sigma_s^2}\, x_s
+ \frac{\alpha_t\sigma_{s|t}^2}{\sigma_s^2}\, x_{\phi^\times}(x_s,s)
}_{\mu_{\phi^\times}(x_s,t,s)}
+ \sigma_{s|t}\frac{\sigma_t}{\sigma_s}\,\epsilon_s,
\qquad \epsilon_s\sim\mathcal{N}(0,I). \quad (4.5.8)
\]


\end{document}
